\section{Cone pictures and torsion-freeness}\label{sec:topology}

Fix graphs supplied by \Cref{prop:no-arrangements}. Attach to the rose \(F\), for each connected component \(\Lambda\) of \(\Gamma\), its abstract cone \(K_\Lambda\) along the labelled map \(\Lambda\to F\). Denote the resulting finite two-dimensional CW complex by \(X\), and put \(G=\pi_1(X)\), based at the rose vertex. In particular, labels of all closed walks in \(\Gamma\) are trivial in \(G\).

\begin{proposition}\label{prop:topology}
The complex \(X\) satisfies \(\pi_2(X)=0\). Moreover, if a path in \(\Gamma\) starts at \(x_A\) or \(x_B\) and ends at a different vertex, its label is nontrivial in \(G\).
\end{proposition}

We prove the two conclusions together, using minimal surface pictures. The proof keeps track of the actual boundary lifts in \(\Gamma\). Removing overlaps with coincident graph lifts is also central to the graphical-diagram arguments of Gruber and Bereznyuk~\cite{Gruber2015,Bereznyuk2019}. Their hypotheses are not imported here: we give the cone-picture reduction needed for \Cref{prop:no-arrangements} directly, without a graphical small-cancellation condition.

\subsection{Representing maps by pictures}

A \emph{cone picture} on a sphere or disk consists of finitely many disjoint tame disks in the surface interior, called \emph{inner disks}, with the following mapping data. Each inner disk has a specified boundary lift to a graph component \(\Lambda\), its filling factors through \(K_\Lambda\), and the complement of the disk interiors maps to \(F\). The boundary lifts are combinatorial closed paths. For a disk domain the outer boundary additionally reads the label of a specified interval path in \(\Gamma\).

\begin{lemma}\label{lem:pictures}
Every essential sphere map to \(X\) can be replaced by an essential sphere map admitting a cone picture. Every null-homotopy of a graph-path label can be represented by a disk cone picture with that boundary label.
\end{lemma}
\begin{proof}
For each component \(\Lambda\), choose based edge loops representing a free basis of \(\pi_1(\Lambda)\). Temporarily replace \(K_\Lambda\) by the space \(Y_\Lambda\) obtained by attaching one disk to \(\Lambda\) along each basis loop. The space \(Y_\Lambda\) is simply connected. The classes of its attaching loops form a basis of \(H_1(\Lambda;\mathbb Z)\), so its cellular chains show that its positive-dimensional homology vanishes. A simply connected acyclic CW complex is contractible, by the Hurewicz and Whitehead theorems; see~\cite[Theorem~4.32 and Corollary~4.33]{Hatcher2002}.

Both \(Y_\Lambda\) and \(K_\Lambda\) contain \(\Lambda\) as a subcomplex. Their contractibility gives maps between them restricting to the identity on \(\Lambda\), and homotopies of the composites to the identities relative to \(\Lambda\): extend the prescribed maps and homotopies over relative cells, using the vanishing of the homotopy groups of the targets. These relative maps and homotopies glue with the identity on \(F\). This works even though \(\Lambda\to F\) is not injective. We obtain an ordinary rose-with-relator-disks model homotopy equivalent to \(X\) relative to \(F\).

That finite model can be triangulated. Take collars inside the relator disks, view their attachments to \(F\) as mapping cylinders, subdivide the collar rectangles so that their bottom edges map injectively to small rose edges, and triangulate. A surface map can be approximated piecewise linearly, preserving its boundary in the graph. If needed, a graph collar homotopy restores the original boundary word in the disk case.

Choose a point in the interior of each relator disk outside the images of degenerate or lower-dimensional simplices. It has finitely many preimages, with disjoint disk neighborhoods that map homeomorphically onto a small disk around the chosen point. Choose these target disks small enough that there are no other preimages. A radial map of each relator disk expands the small target disk to the whole relator disk and sends the remaining annulus to its attaching boundary. This map is homotopic to the identity relative to \(F\).

After composition, the complement of the preimage disks maps to \(F\), and each preimage disk reads an attaching loop, with one of its two orientations. Return to \(X\) using the relative equivalences of the individual cone substitutes. Each disk then factors through its specified abstract cone. Essentiality of a sphere, or the prescribed disk-boundary word, is preserved.
\end{proof}

\subsection{Minimizing a failure picture}

Suppose that one conclusion of \Cref{prop:topology} fails. In the sphere case, consider all essential sphere pictures; the homotopy class is allowed to vary. In the disk case, fix two distinct vertices witnessing a failure, with the first one equal to \(x_A\) or \(x_B\), and allow the interval path between them to vary. By \Cref{lem:pictures}, the appropriate class of pictures is nonempty.

Choose a picture of least total boundary length, counting all inner boundaries and, in the disk case, the outer path. This length is a nonnegative integer. Cyclically tighten every inner loop in its own graph component. The tightening can be performed on a collar, with the smaller disk filled through the same cone. The old filling and this replacement are homotopic relative to the old boundary, since both factor through that contractible cone. A constant boundary disk can be removed. Thus every inner boundary is nonempty and cyclically immersed.

Tighten the outer interval relative to its actual graph endpoints. It remains nonempty because those endpoints are distinct. There must be at least one inner disk: a graph has no essential sphere, and a nonempty immersed graph path has a nonempty freely reduced word under the immersion to the rose, so its label cannot be trivial in \(\pi_1(F)\).

For a disk-domain picture, cap the exterior conceptually by an extra disk. Its boundary carries the interval path, with a single marked break from the endpoint back to the initial vertex. No map of that extra disk into \(X\) is required. Orient the sphere so that the coherent exterior-disk boundary orientation reads this path from its prescribed root to its endpoint. We now have an oriented sphere with ordinary inner disks and at most one exceptional exterior disk.

In the surface region mapping to \(F\), take transverse preimages of one interior point in each rose edge. This can again be arranged by PL approximation relative to the boundary words, using collars to straighten their parametrizations. The preimages are disjoint circles and arcs. Ignore the circles. Every letter occurrence has exactly one arc endpoint, giving a complete pairing. Transverse orientation over a rose edge gives opposite crossing signs at the two ends of an arc along the oriented boundary of the complement. Switching to coherent disk-boundary orientations reverses both signs. Consequently the paired letters are inverse in those orientations.

It remains to establish the reduction condition \eqref{eq:edge-reduction}. The following surgery is the reason for minimizing over boundary length and, for spheres, over all nonzero classes.

\subsection{Cancelling a band over the same lifted edge}

Suppose an arc joins occurrences of the same underlying unoriented edge in \(\Gamma\). A sufficiently narrow band around that arc maps into the corresponding open rose edge. It lifts to the specified common open graph edge, because the map of that open edge to the rose is a homeomorphism. The lift agrees with the specified lifts at both attachment intervals. In particular, any incident inner disks and the band together factor through one common abstract cone.

Join the incident disks along the band. If it returns to one disk, the result is an annulus in the sphere. Boundary paths follow the original paths cut at the two occurrences and spliced along the band. They lift to the same graph component, except for the one retained marked break when the exterior is involved.

Both spliced joins contract along the common graph edge. To check this concretely, orient it as \(e:u\to v\). The two matched occurrences are \(e\) and \(e^{-1}\). The local possibilities are:
\begin{center}
\renewcommand{\arraystretch}{1.18}
\begin{tabular}{@{}p{0.24\textwidth}p{0.33\textwidth}p{0.33\textwidth}@{}}
\toprule
Incident disks & Old lifted paths & New lifted paths\\
\midrule
Two inner disks & \(ea,\ e^{-1}b\) & \(ab\)\\
Exterior and inner & \(aec:x\to y,\ e^{-1}b\) & \(abc:x\to y\)\\
One inner disk twice & \(eae^{-1}b\) & \(a,\ b\)\\
Exterior twice & \(aebe^{-1}c:x\to y\) & \(ac:x\to y\)\\
\bottomrule
\end{tabular}
\end{center}
The symbols denote graph paths with the endpoints forced by the displayed concatenations. For example, in the last row \(a:x\to u\), \(b:v\to v\), and \(c:u\to y\). The other new boundary is \(b\), up to orientation. Thus the break-containing path \(ac\) has exactly the original distinct endpoints \(x,y\).

Geometrically, each band side joins two residual edge fragments that go from one endpoint of \(e\) back to that same endpoint through its interior. These fragments contract relative to their endpoints. Taken together, the two contractions remove the two complete edge occurrences and add no letters. Neither contraction passes through the marked break. Total new boundary length is smaller by two before any further tightening.

Here are the resulting modifications of the entire picture.
\begin{enumerate}
\item If the band joins two distinct inner disks, replace their disk-and-band union by one cone disk. The replacement is homotopic relative to its boundary because the entire map there factors through the common contractible cone. Then tighten its boundary.
\item If the band joins an inner disk to the exterior disk, replace the exterior by their joined disk. Keep the map on its complement and tighten the new outer path relative to the same endpoints.
\item If both ends meet one inner disk, remove its annulus with the band and cap the boundaries on the two complementary parts using that same cone. For a disk-domain failure, retain the part containing the exterior. For a sphere failure, at least one capped sphere remains essential, as proved below.
\item If both ends meet the exterior disk, the resulting annulus has the marked break on one boundary component. Retain as the new domain the complementary disk adjacent to that boundary. The last row of the table shows that its outer path has the same graph endpoints.
\end{enumerate}

\begin{figure}[htbp]
\centering
\begin{tikzpicture}[x=0.83cm,y=0.83cm,>=Latex,font=\small,line width=0.6pt]
\begin{scope}
 \fill[gray!17,even odd rule] (0,0) circle (1.13) (0,0) circle (0.52);
 \draw (0,0) circle (1.13);
 \draw (0,0) circle (0.52);
 \node at (0,0) {$a$};
 \node at (-1.47,0) {$b$};
 \node[align=center] at (0,-1.65) {disk and band\\(an annulus)};
\end{scope}
\draw[->] (1.65,0)--(3.7,0);
\node[align=center,font=\scriptsize,anchor=south] at (2.68,0.16) {remove\\annulus;\\cap\\complements};
\begin{scope}[shift={(5.05,0)}]
 \fill[gray!5] (0,0) circle (1.03);
 \draw (0,0) circle (1.03);
 \fill[gray!17] (0,-0.38) ellipse (0.83 and 0.25);
 \draw (0,-0.38) ellipse (0.83 and 0.25);
 \node at (0,-0.38) {$a$};
 \node at (0,0.32) {$S_1$};
\end{scope}
\node at (6.53,0) {$+$};
\begin{scope}[shift={(8,0)}]
 \fill[gray!5] (0,0) circle (1.03);
 \draw (0,0) circle (1.03);
 \fill[gray!17] (0,-0.38) ellipse (0.83 and 0.25);
 \draw (0,-0.38) ellipse (0.83 and 0.25);
 \node at (0,-0.38) {$b$};
 \node at (0,0.32) {$S_2$};
\end{scope}
\node at (6.5,-1.65) {$[S]=[S_1]+[S_2]\quad\text{in }H_2(\widetilde X;\mathbb Z)$};
\end{tikzpicture}


\caption{An inner disk joined to itself by a band forms an annulus. Cutting away the annulus leaves two complementary disks. Capping their new boundaries through the same abstract cone gives two spheres whose homology classes sum to the original class in the universal cover. The labels \(a,b\) indicate the shortened boundary paths; the picture is schematic.}
\label{fig:band-surgery}
\end{figure}

For the essentiality assertion in the third case, illustrated in \Cref{fig:band-surgery}, lift the original sphere map to the universal cover \(\widetilde X\). The map \(K_\Lambda\to X\) has a lift because the cone is simply connected. Choose this lift to agree with the lifted sphere at one point of the annulus. Since the annulus is connected, uniqueness of lifts gives agreement throughout it. Fill both new caps in \(K_\Lambda\) and use this same lift.

With compatible orientations, the original spherical class in \(H_2(\widetilde X;\mathbb Z)\) is the sum of the two capped classes. The difference between the original class and this sum is represented by the annulus together with the reverse caps, all mapping through the contractible cone, so it is zero in homology. Since \(\widetilde X\) is simply connected, the Hurewicz map sending a sphere to its fundamental homology class identifies \(\pi_2(\widetilde X)\) with \(H_2(\widetilde X;\mathbb Z)\)~\cite[Theorem~4.37]{Hatcher2002}. A nonzero original class therefore leaves a nonzero capped class. This uses no prior vanishing of \(\pi_2(X)\).

In every case the retained surface still witnesses the chosen kind of failure, and its total boundary length has strictly decreased. For the split sphere this follows because the total new length over both parts is already smaller; keeping just one part cannot increase it. Collar homotopies restore combinatorial boundaries after the contractions, and further tightening only decreases length. In the exterior cases, the new graph path has the same fixed distinct endpoints. This contradicts the chosen minimum.

Thus every paired occurrence uses a different underlying graph edge from its partner. The minimized picture is now an arrangement of \Cref{def:arrangement}: it has at least one ordinary cyclic boundary, at most one rooted exceptional interval with one break, a full planar inverse-letter pairing, and the reduction condition. \Cref{prop:no-arrangements} excludes it. This proves \Cref{prop:topology}.

\subsection{A finite-dimensional obstruction to torsion}

\begin{corollary}\label{cor:torsion-free}
The universal cover of \(X\) is contractible, and \(G\) is torsion-free.
\end{corollary}
\begin{proof}
The simply connected two-dimensional cover \(\widetilde X\) has \(H_2=0\) by \Cref{prop:topology} and Hurewicz, has \(H_1=0\), and has no cellular homology above dimension two. It is therefore contractible by the homological form of Whitehead's theorem~\cite[Corollary~4.33]{Hatcher2002}.

For completeness, the torsion obstruction can be seen directly from resolutions; compare~\cite[Chapter~VIII, Section~2]{Brown1982}. The augmented cellular chains of \(\widetilde X\) give an exact free \(\mathbb ZG\)-resolution of \(\mathbb Z\) of length two. If \(G\) had a nontrivial finite-order element, a power would generate a cyclic subgroup \(H=\langle g\rangle\) of prime order \(\ell\). Restricting the resolution to \(\mathbb ZH\) preserves freeness, since \(\mathbb ZG\) is free as a left \(\mathbb ZH\)-module on coset representatives.

The cyclic group has the exact periodic resolution
\[
\cdots\longrightarrow\mathbb ZH
 \xrightarrow{\,N\,}\mathbb ZH
 \xrightarrow{\,g-1\,}\mathbb ZH
 \xrightarrow{\,N\,}\mathbb ZH
 \xrightarrow{\,g-1\,}\mathbb ZH
 \longrightarrow\mathbb Z\longrightarrow0,
 \qquad N=1+g+\cdots+g^{\ell-1}.
\]
Exactness follows by comparing the coefficients of \(1,g,\ldots,g^{\ell-1}\): the kernel of \(g-1\) consists of the multiples of \(N\), and the kernel of \(N\) is the augmentation ideal generated by \(g-1\). With trivial coefficients \(\mathbb F_\ell\), all differentials in the corresponding cochain complex vanish. Hence
\(H^k(H;\mathbb F_\ell)\cong\mathbb F_\ell\) in every degree. The comparison theorem for projective resolutions makes this cohomology independent of the chosen resolution~\cite[Chapter~I, Theorem~7.5, p.~24]{Brown1982}. This contradicts the vanishing above degree two forced by the restricted finite-length free resolution.
\end{proof}
